\documentclass[10pt,a4paper]{amsart}
\usepackage[margin=19mm]{geometry}
\usepackage{amsmath,amssymb,mathtools}
\usepackage[scr=rsfs,cal=euler]{mathalfa}
\usepackage{xfrac}
\usepackage[unicode,hidelinks,bookmarks=false]{hyperref}
\newcommand{\ie}{i.e.\ }
\newcommand{\cf}{cf.\ }
\newcommand{\resp}{respectively\ }
\newcommand{\Subsection}{\subsection}
\providecommand{\nobreakdash}{\nobreak\hbox{-}}
\setlength{\emergencystretch}{2em}
\multlinegap0pt
\usepackage{amssymb}
\newtheorem{theorem}{Theorem}
\newtheorem{lemma}{Lemma}
\title{The capillary $L_p$ dual Minkowski problem for $p>q$ in higher dimensions}
\author{Ya Gao$^{\dagger,\ast}$}
\date{Source: arXiv:2609.07158v1 (CC BY 4.0)}
\begin{document}
\maketitle

\noindent\emph{Source and licence.} The capillary $L_p$ dual Minkowski problem for $p>q$ in higher dimensions. Version arXiv:2609.07158v1 (CC BY 4.0), distributed by arXiv under the Creative Commons Attribution 4.0 International licence, http://creativecommons.org/licenses/by/4.0/, as recorded in the arXiv licence metadata for this exact version and shown on the abstract page https://arxiv.org/abs/2609.07158v1. Full licence notice: \texttt{licenses/arxiv-2609.07158-CC-BY-4.0.txt}.\par\smallskip

\noindent\textit{Editorial scope.} Counted unit: the article's theorem estimate (source0.tex lines 602--615). The article's complete original proof of it is delivered with it (source0.tex lines 617--1088, one \texttt{proof} environment of 472 lines). No mathematics is written, repaired or re-derived: the statement and the proof are the article's own lines, in the article's own notation, and no retained line is substituted or re-flowed.  Retained uncounted as the source-local context the proof needs: lines 258--265.  Nothing was omitted from any retained span, so the package carries no omission record.  No retained line contains a citation, so the wrapper carries no bibliography and no bibliography entry.  No retained line contains a figure environment, an image-inclusion command or drawing code, so no image is delivered and the package declares no asset dependency.  Editorial formatting only, in addition: an AMS \texttt{amsart} preamble that restates the article's own declarations and macros used by the retained lines, namely the environments \texttt{theorem}, \texttt{lemma} on separate counters, together with the standard packages those lines need.  The retained environments print in this document as Theorem 1, Lemma 1 in order of appearance, whereas the article prints the same items with the numbering of its own full document.  The complete native source of the article as distributed in the arXiv e-print for this exact version is bundled unchanged as \texttt{sources/209605/source0.tex}.
\begin{equation}\label{Eq}
\left\{
\begin{aligned}
&\det\left(\nabla^2 h + h\sigma\right)=fh^{p-1}\left(h^2 + |\nabla h|^2\right)^{\frac{n+1-q}{2}}, ~~&&in~~C_{\theta}\\
&\nabla_{\mu} h=\cot\theta\cdot h,  ~~&& on~~\partial C_{\theta}
\end{aligned}
\right.
\end{equation}

\begin{theorem}[{\bf{$C^2$ estimate}}]\label{c2 estimate}
Let $n\geq 3$, $p>q$, and $\theta\in (0,\frac{\pi}{2})$. Suppose that $h$ is a
positive $C^4$ strictly convex solution to Eq. \eqref{Eq}. Assume that
\[
        0<m\leq h\leq M,\qquad |\nabla h|\leq K,
\]
then there holds 
\begin{equation}\label{d1}
        \max_{C_{\theta}}|\nabla^2h|\leq C,
\end{equation}
where the positive constant $C$ depending only on
$n,p,q,\theta,m,M,K,\min_{C_{\theta}}f$ and
$\Vert f\Vert_{C^2(C_{\theta})}$, 
\end{theorem}

\begin{proof}
The equation and the lower-order bounds give determinant bounds
\[
        D_-\leq\det A\leq D_+,
\]
where
\[
        D_- =
        \min_{C_{\theta}}f\cdot\min\{m^{p-1},M^{p-1}\}
        \cdot\min\{(m^2)^{\beta/2},(M^2+K^2)^{\beta/2}\}
\]
and
\[
        D_+ =
        \max_{C_{\theta}}f\cdot\max\{m^{p-1},M^{p-1}\}
        \cdot\max\{(m^2)^{\beta/2},(M^2+K^2)^{\beta/2}\}.
\]
Let
\begin{equation}\label{auxiliary Q}
        Q:=\log\sigma_1(A)+ah+b|\nabla h|^2+\log\sigma_{n-1}(A),
\end{equation}
where $b>0$ and $a<0$ will be chosen later. Suppose that $Q$ attains its maximum value at some point $\xi_0\in C_{\theta}$. We divide this proof into two cases: either $\xi_0\in \partial C_{\theta}$ or $\xi_0\in C_{\theta} \backslash \partial C_{\theta}$.

\noindent\textbf{Case 1. $\xi_0\in C_{\theta} \backslash \partial C_{\theta}$.}
In this case, we choose an orthonormal frame $\{e_i\}_{i=1}^n$ around $\xi_0$, such that $A_{ij}=(h_{ij}+h\delta_{ij})$ is diagonal, hence $F^{ij}:=\frac{\partial F(A)}{\partial A_{ij}}$ and $h_{ij}$ are also diagonal at $\xi_0$. In this case, we have
\[
        A_{ij,k}=A_{ik,j},\qquad
        A_{ii,jj}=A_{jj,ii}+A_{ii}-A_{jj}.
\]

Since
\[
        \sigma_{n-1}(A)=\det(A)\,S,
\]
define
\[
        W_i=\frac{1}{\sigma_1(A)}+\lambda_i^{-1}
        -\frac{\lambda_i^{-2}}{S}.
\]
Then at $\xi_0$, we have
\begin{equation}\label{Q stationary}
        0=Q_k
        =\sum_iW_iA_{ii,k}+\bigl(a+2b(\lambda_k-h)\bigr)h_k.
\end{equation}

We need the quadratic third-derivative terms that occur after contracting
second derivatives with $A^{ij}$.

\begin{lemma}\label{third forms}
At a diagonal point where $A_{ij,k}$ is symmetric in all three indices, set
\[
\begin{split}
        \mathcal G_{\sigma_1}
        &=\frac{1}{\sigma_1(A)}
        \sum_{i,j,k}\lambda_i^{-1}\lambda_j^{-1}A_{ij,k}^2
        -\frac{1}{\sigma_1(A)^2}
        \sum_k\lambda_k^{-1}\left(\sum_iA_{ii,k}\right)^2,
\end{split}
\]
and
\[
\begin{split}
        \mathcal G_S
        &=\frac{1}{S}
        \sum_{i,j,k}
        \left(2\lambda_k^{-1}\lambda_i^{-2}\lambda_j^{-1}
        -\lambda_k^{-2}\lambda_i^{-1}\lambda_j^{-1}\right)A_{ij,k}^2\\
        &\qquad
        -\frac{1}{S^2}
        \sum_k\lambda_k^{-1}
        \left(\sum_i\lambda_i^{-2}A_{ii,k}\right)^2.
\end{split}
\]
Then $\mathcal G_{\sigma_1}\geq0$ and $\mathcal G_S\geq0$.
\end{lemma}

\noindent\emph{Proof of Lemma \ref{third forms}.}
Fix the derivative index $k$ and denote $v_i=A_{ii,k}$. For
$\sigma_1(A)\cdot\mathcal G_{\sigma_1}$, the diagonal coefficients after grouping all terms
forced by the three-index symmetry are
\[
        d_k=\lambda_k^{-2},\qquad
        d_i=\lambda_i^{-2}+2\lambda_i^{-1}\lambda_k^{-1}\quad (i\ne k),
\]
and the negative term is
\[
        -\frac{\lambda_k^{-1}}{\sigma_1(A)}
        \left(\sum_i v_i\right)^2.
\]
Cauchy-Schwarz inequality gives
\[
        \left(\sum_i v_i\right)^2
        \leq \left(\sum_i d_i v_i^2\right)\left(\sum_i d_i^{-1}\right).
\]
Here $d_k^{-1}=\lambda_k^2$, and for $i\ne k$,
\[
        d_i^{-1}
        =\frac{\lambda_i^2\lambda_k}{\lambda_k+2\lambda_i}
        \leq \lambda_i\lambda_k.
\]
Thus
\[
        \sum_i d_i^{-1}\leq \lambda_k\sum_i\lambda_i
        =\lambda_k\sigma_1(A),
\]
which makes the grouped diagonal part nonnegative. If $i,j,k$ are pairwise
distinct, the total coefficient of the six symmetric permutations is
\[
        2(\lambda_i^{-1}\lambda_j^{-1}
        +\lambda_i^{-1}\lambda_k^{-1}
        +\lambda_j^{-1}\lambda_k^{-1})\geq0.
\]
Therefore $\mathcal G_{\sigma_1}\geq0$.

Similarly, for $S\cdot\mathcal G_S$, again fix $k$ and write $v_i=A_{ii,k}$. The diagonal
coefficients are
\[
        d_k=\lambda_k^{-4},\qquad
        d_i=\lambda_i^{-2}\lambda_k^{-1}
        (2\lambda_i^{-1}+\lambda_k^{-1})\quad (i\ne k),
\]
and the negative term is
\[
        -\frac{\lambda_k^{-1}}{S}
        \left(\sum_i\lambda_i^{-2}v_i\right)^2.
\]
Cauchy-Schwarz inequality gives
\[
        \left(\sum_i\lambda_i^{-2}v_i\right)^2
        \leq \left(\sum_i d_i v_i^2\right)
        \left(\sum_i\frac{\lambda_i^{-4}}{d_i}\right).
\]
It is enough to prove
\[
        \lambda_k^{-1}
        +\sum_{i\ne k}
        \frac{\lambda_i^{-2}}{2\lambda_i^{-1}+\lambda_k^{-1}}
        \leq S,
\]
and this follows from
\[
        \frac{\lambda_i^{-2}}{2\lambda_i^{-1}+\lambda_k^{-1}}
        \leq \lambda_i^{-1}.
\]
For pairwise distinct $i,j,k$, the coefficient of the six symmetric
permutations is
\[
        2\lambda_i^{-1}\lambda_j^{-1}\lambda_k^{-1}
        (\lambda_i^{-1}+\lambda_j^{-1}+\lambda_k^{-1})\geq0.
\]
Hence $\mathcal G_S\geq0$. This proves the local lemma. \hfill$\Box$

Differentiating $\log\det A=\tilde f$ twice we have
\begin{equation}\label{s4-1}
        \tilde f_{ii}
        =\sum_k\lambda_k^{-1}A_{kk,ii}
        -\sum_{k,s}\lambda_k^{-1}\lambda_s^{-1}A_{ks,i}^2,
\end{equation}
so, for $\log\sigma_1(A)$ we know
\[
        \sum_i\lambda_i^{-1}(\log\sigma_1(A))_{ii}
        =\mathcal G_{\sigma_1}+\frac{1}{\sigma_1(A)}\sum_i\tilde f_{ii}
        +S-\frac{n^2}{\sigma_1(A)}.
\]
For the inverse term $S$, by directly calculate
\[
        S_k=-\sum_i\lambda_i^{-2}A_{ii,k},
\]
and
\begin{equation}\label{s4-2}
        S_{kk}
        =-\sum_i\lambda_i^{-2}A_{ii,kk}
        +2\sum_{i,j}\lambda_i^{-2}\lambda_j^{-1}A_{ij,k}^2.
\end{equation}
Using the formula $A_{ii,jj}=A_{jj,ii}+A_{ii}-A_{jj}$ and \eqref{s4-1}, we have
\[
        \sum_k\lambda_k^{-1}A_{ii,kk}
        =\tilde f_{ii}
        +\sum_{k,s}\lambda_k^{-1}\lambda_s^{-1}A_{ks,i}^2
        +\lambda_i\cdot S-n.
\]
Then the formula \eqref{s4-2} can be written as
\[
\begin{split}
        \sum_k\lambda_k^{-1}S_{kk}
        &=-\sum_i\lambda_i^{-2}\tilde f_{ii}
        -S^2+n\sum_i\lambda_i^{-2}\\
        &\qquad
        +\sum_{i,j,k}
        \left(2\lambda_k^{-1}\lambda_i^{-2}\lambda_j^{-1}
        -\lambda_k^{-2}\lambda_i^{-1}\lambda_j^{-1}\right)A_{ij,k}^2.
\end{split}
\]
After subtracting $S^{-2}\sum_k\lambda_k^{-1}
(\sum_i\lambda_i^{-2}A_{ii,k})^2$, this yields
\begin{equation}\label{contracted log S}
        \sum_k\lambda_k^{-1}(\log S)_{kk}
        =\mathcal G_S-\frac{1}{S}\sum_i\lambda_i^{-2}\tilde f_{ii}
        -S+\frac{n\sum_i\lambda_i^{-2}}{S}.
\end{equation}
Also
\[
        \sum_i\lambda_i^{-1}h_{ii}=n-hS
\]
and
\[
\begin{split}
        \sum_i\lambda_i^{-1}(|\nabla h|^2)_{ii}
        &=2\sum_kh_k\tilde f_k+2\sigma_1(A)-4nh
        +2h^2S-2\sum_i\lambda_i^{-1}h_i^2.
\end{split}
\]
Combining the preceding formulas with $\log\sigma_{n-1}(A)=\tilde f+\log S$, at $\xi_0$,
we get
\begin{equation}\label{contracted Q}
\begin{split}
        0&\geq\sum_i\lambda_i^{-1}Q_{ii}\\
        &=\mathcal G_{\sigma_1}+\mathcal G_S+\sum_iW_i\tilde f_{ii}
        -\frac{n^2}{\sigma_1(A)}+a(n-hS)\\
        &\qquad
        +b\left(2\sum_kh_k\tilde f_k+2\sigma_1(A)-4nh
        +2h^2S-2\sum_i\lambda_i^{-1}h_i^2\right)
        +\frac{n\sum_i\lambda_i^{-2}}{S}.
\end{split}
\end{equation}
Lemma \ref{third forms} makes the first two terms nonnegative, and easy find that the last term is also
nonnegative.

Since $\tilde f:=(p-1)\log h+\frac{\beta}{2}\log P+\log f$, by directly calculate we know
\begin{equation}\label{f second derivatives}
\begin{split}
        \tilde f_{ii}
        &=\beta\left(
        \frac{\lambda_i^2}{P}
        -\frac{2h_i^2\lambda_i^2}{P^2}
        +\frac{\sum_k h_kA_{ii,k}}{P}
        -\frac{h\lambda_i}{P}\right)\\
        &\qquad
        +(p-1)\left(\frac{\lambda_i}{h}-1-\frac{h_i^2}{h^2}\right)
        +(\log f)_{ii}.
\end{split}
\end{equation}
Using $Q_k=0$, then
\[
      \frac{\beta}{P}
        \sum_kh_k\sum_iW_iA_{ii,k}=  -\frac{a\beta|\nabla h|^2}{P}
        -\frac{2b\beta}{P}\sum_kh_k^2\lambda_k
        +\frac{2b\beta h|\nabla h|^2}{P}.
\]
On the other hand,
\[
        2b\sum_kh_k\tilde f_k
        =2b(p-1)\frac{|\nabla h|^2}{h}
        +\frac{2b\beta}{P}\sum_kh_k^2\lambda_k
        +2b\sum_kh_k(\log f)_k.
\]
The terms containing $\sum_kh_k^2\lambda_k$ cancel.

Since
$P\geq m^2$ and $|h_i|^2\leq K^2$, the $\lambda_i^2$ terms above
are bounded from below by $-C\sum_iW_i\lambda_i^2$. For $W_i$ we have
\[
        0\leq W_i,\qquad
        \sum_iW_i\lambda_i^2\leq 2\sigma_1(A),\qquad
        \sum_iW_i\lambda_i= n,\qquad
        \sum_iW_i\leq C+S,
\]
where we used $\sigma_1(A)\geq nD_-^{1/n}$. Thus the terms containing
$\beta$ and $p-1$ are controlled for all signs of these numbers. The
$(\log f)_{ii}$ terms are bounded below by some negative constant $-C$, and
\[
        2h^2S-2\sum_i\lambda_i^{-1}h_i^2\geq -2K^2S.
\]

So the inequality \eqref{contracted Q} can be written as
\begin{equation}\label{interior inequality}
\begin{split}
        0\geq
        (2b-C)\sigma_1(A)+(-am-2bK^2-C)S
        -C(1+|a|+b).
\end{split}
\end{equation}
Choose $b=b_1>C+1$. Then choose $a<a_1<0$ such that
\[
        -a_1m-2bK^2-C\geq 1.
\]
Then at $\xi_0$, the formula \eqref{interior inequality} can be written as 
\[
        0\geq \sigma_1(A) - C(1+|a|+b).
\]
So we conclude that $\sigma_1(\xi_0)\leq C$, where the positive constant $C$ depends on $n, p, q, \min_{C_{\theta}}f, \min_{C_{\theta}}h, \Vert f\Vert_{C^2 (C_{\theta})}$ and $\Vert h\Vert _{C^0 (C_{\theta})}$.

\noindent\textbf{Case 2. $\xi_0\in \partial C_{\theta}$}
We choose an orthonormal frame $\{e_i\}_{i=1}^n$ around $\xi_0\in\partial C_{\theta}$ satisfying $e_n=\mu$ at $\xi_0$. So we know 
\[
        \lambda_n=A_{nn},\qquad
        \lambda_{\alpha}=A_{\alpha\alpha}, \quad \alpha\in 1,2,\cdots, n-1
\]

At the boundary, with Robin condition $h_n=\cot\theta\,h$, we know that $h_{\alpha n}=0$.
By directly calculate, we know the normal derivative of $\tilde{f}$ is 
\[
        \tilde f_n
        =(p-1)\cot\theta
        +\beta\cot\theta\,
        \frac{h\lambda_n}{P}
        +(\log f)_n,
\]
by $P_n=2\cot\theta\,h\lambda_n$. Differentiating
$\log\det A=\tilde f$ in the normal direction and using
$S=\lambda_n^{-1}+\sum_{\alpha<n}\lambda_{\alpha}^{-1}$ we have
\[
        (\lambda_n)_n
        =\lambda_n\tilde f_n
        -\cot\theta\,\lambda_n(\lambda_nS-n).
\]
Then
\[
\begin{split}
        (\sigma_1(A))_n
        &=(\lambda_n)_n+\sum_{\alpha<n}(\lambda_{\alpha})_n\\
        &=\lambda_n\tilde f_n
        +\cot\theta\bigl(2n\lambda_n-\sigma_1(A)-\lambda_n^2S\bigr),
\end{split}
\]
and
\[
\begin{split}
        S_n
        &=-\lambda_n^{-2}(\lambda_n)_n
        -\sum_{\alpha<n}\lambda_{\alpha}^{-2}(\lambda_{\alpha})_n\\
        &=-\frac{\tilde f_n}{\lambda_n}
        -\cot\theta\,\lambda_n
        \sum_{\alpha<n}\left(\frac1{\lambda_{\alpha}}
        -\frac1{\lambda_n}\right)^2.
\end{split}
\]
It follows from $\sigma_{n-1}(A)=\det(A)S$ that
\[
        (\log\sigma_{n-1}(A))_n
        =\left(1-\frac1{\lambda_nS}\right)\tilde f_n
        -\frac{\cot\theta\,\lambda_n}{S}
        \sum_{\alpha<n}\left(\frac1{\lambda_{\alpha}}
        -\frac1{\lambda_n}\right)^2.
\]
Together with $|\nabla h|^2_n=2\cot\theta\,h(\lambda_n-h)$ and
\[
        2n-\frac{\sigma_1(A)}{\lambda_n}-\lambda_nS
        =-\sum_{\alpha<n}\left(\frac{\lambda_n}{\lambda_{\alpha}}
        +\frac{\lambda_{\alpha}}{\lambda_n}-2\right). 
\]
Then, at $\xi_0$, we have
\begin{equation}\label{boundary Q derivative}
\begin{split}
        0\leq \frac{\sigma_1(A)}{\lambda_n}Q_n
        &=\left(1+\frac{\sigma_1(A)}{\lambda_n}
        -\frac{\sigma_1(A)}{\lambda_n^2S}\right)\tilde f_n\\
        &\qquad
        -\cot\theta
        \sum_{\alpha<n}\left(\frac{\lambda_n}{\lambda_{\alpha}}
        +\frac{\lambda_{\alpha}}{\lambda_n}-2\right)\\
        &\qquad
        +a\cot\theta\,h\frac{\sigma_1(A)}{\lambda_n}
        +2b\cot\theta\,h(\lambda_n-h)\frac{\sigma_1(A)}{\lambda_n}\\
        &\qquad
        -\cot\theta\,\frac{\sigma_1(A)}{S}
        \sum_{\alpha<n}\left(\frac1{\lambda_{\alpha}}
        -\frac1{\lambda_n}\right)^2.
\end{split}
\end{equation}

First suppose $\lambda_n\leq m/2$. Then $\lambda_n-h\leq -m/2$, so the
 term
 $$2b\cot\theta h(\lambda_n -h)\frac{\sigma_1(A)}{\lambda_n}$$
 is nonpositive. Then we can written \eqref{boundary Q derivative} as

 \begin{equation}\label{boundary Q derivative 2}
\begin{split}
        0&\leq \left(1+\frac{\sigma_1(A)}{\lambda_n}
        -\frac{\sigma_1(A)}{\lambda_n^2S}\right)\tilde f_n
        -\cot\theta
        \sum_{\alpha<n}\left(\frac{\lambda_n}{\lambda_{\alpha}}
        +\frac{\lambda_{\alpha}}{\lambda_n}-2\right)\\
        &\qquad
       +a\cot\theta\,h\frac{\sigma_1(A)}{\lambda_n} -\cot\theta\,\frac{\sigma_1(A)}{S}
        \sum_{\alpha<n}\left(\frac1{\lambda_{\alpha}}
        -\frac1{\lambda_n}\right)^2  \\ 
        & 
        \leq \left(1+\frac{\sigma_1(A)}{\lambda_n}
        -\frac{\sigma_1(A)}{\lambda_n^2S}\right)\tilde f_n
        -\cot\theta
        \sum_{\alpha<n}\left(\frac{\lambda_n}{\lambda_{\alpha}}
        +\frac{\lambda_{\alpha}}{\lambda_n}-2\right) +a\cot\theta\,h\frac{\sigma_1(A)}{\lambda_n}\\ 
        &\leq 
        \left(1+\frac{\sigma_1(A)}{\lambda_n}
        -\frac{\sigma_1(A)}{\lambda_n^2S}\right)\tilde f_n +a\cot\theta\,h\frac{\sigma_1(A)}{\lambda_n},
\end{split}
\end{equation}
 The last inequality holds by using $z+z^{-1}-2 = \frac{(z-1)^2}{z} ~~(z>0)$.
We know that 
$$ \tilde f_n
        =(p-1)\cot\theta
        +\beta\cot\theta\,
        \frac{h\lambda_n}{P}
        +(\log f)_n \leq C+C\lambda_n,
        $$
where the positive constant $C$ depend on $n, p, q, \min_{C_{\theta}}f, \min_{C_{\theta}}h, \Vert f\Vert_{C^2 (C_{\theta})}$ and $\Vert h\Vert _{C^0 (C_{\theta})}$. Since we suppose $\lambda_n \leq m/2$, so we have $\tilde{f}_n\leq C$. So we can rewritten the formula \eqref{boundary Q derivative 2} as 
\begin{equation}\label{s4-3}
0\leq C+\left(C+a\cot\theta\,m\right)
        \frac{\sigma_1(A)}{\lambda_n}.
\end{equation}

Then choose $a<a_2<0$ such that $C+a_2 \cot\theta m<-1$. Then we have
$$0\leq C-\frac{\sigma_1(A)}{\lambda_n}\leq C-\frac{2}{m}\sigma_1(A), $$
i.e.,
$$\sigma_1(A)\leq \frac{mC}{2}.$$

Next we consider the case that $\lambda_n\geq m/2$. In this case we know 
$\tilde f_n\leq C(1+\lambda_n)$, and
\[
        \left(1+\frac{\sigma_1(A)}{\lambda_n}
        -\frac{\sigma_1(A)}{\lambda_n^2S}\right)\tilde f_n
        \leq C(1+\sigma_1(A)).
\]

Then the formula \eqref{boundary Q derivative} can be written as
\begin{equation}\label{boundary inequality} 
\begin{split}
        0&\leq 
        C(1+\sigma_1(A))
        +a\cot\theta\,h\frac{\sigma_1(A)}{\lambda_n}-\cot\theta\,\frac{\sigma_1(A)}{S}
        \sum_{\alpha<n}\left(\frac1{\lambda_{\alpha}}
        -\frac1{\lambda_n}\right)^2\\
        &\leq 
        C(1+\sigma_1(A))-\cot\theta\,\frac{\sigma_1(A)}{S}
        \sum_{\alpha<n}\left(\frac1{\lambda_{\alpha}}
        -\frac1{\lambda_n}\right)^2.
\end{split}
\end{equation}

By Cauchy-Schwarz inequality,
\[
        \sum_{\alpha<n}\left(\frac1{\lambda_{\alpha}}
        -\frac1{\lambda_n}\right)^2
        \geq\frac1{n-1}\left(S-\frac n{\lambda_n}\right)^2.
\]
If $S\geq4n/m$, then $\lambda_n\geq m/2$ implies
$S-n/\lambda_n\geq S/2$. Therefore we have
\begin{equation}\label{s4-4}
        0\leq C(1+\sigma_1(A))
        -\frac{\cot\theta}{4(n-1)}\sigma_1(A)S.
\end{equation}
Since $\sigma_1(A)\geq nD_-^{1/n}>0$, divide both sides by $\sigma_1(A)$, then we have
$$0\leq C - \frac{\cot\theta}{4(n-1)}S,$$
so $S$ has a positive upper bound. If $S<4n/m$, it is already giving such a bound.
Hence $S\leq C$. Each eigenvalue satisfies
$\lambda_i\geq S^{-1}$, so we have
\[
        \lambda_i
        =\frac{\det A}{\prod_{j\ne i}\lambda_j}
        \leq D_+S^{n-1}\leq C.
\]
Thus $\sigma_1(A)(\xi_0)\leq C$. In conclusion, by setting $a=\min\{a_1, a_2\}, b=b_1$, we can obtain $\sigma_1(A)(\xi_0)\leq C$ at the maximum point $\xi_0\in C_{\theta}$ of $Q$. Then at $\xi_0$, we have 
$$\sigma_{n-1}(A)(\xi_0)\leq n\sigma_1(A)(\xi_0)^{n-1}\leq nC^{n-1},$$
then 
$$Q(\xi_0)\leq C.$$
On the other hand, by MacLaurin inequality, we have
$$\sigma_{n-1}(A)\geq n(\det A)^{(n-1)/n}\geq nD^{(n-1)/n}_-=: c_0>0.$$
In conclusion, at any point $\xi\in C_{\theta}$, we have 
$$\log \sigma_1(A)(\xi) + \log c_0 +aM\leq Q(\xi)\leq Q(\xi_0)\leq C,$$
that means $\sigma_1(A)\leq C$ in $C_{\theta}$.
\end{proof}

\end{document}
